\documentclass{article}
\usepackage[utf8]{inputenc}
\usepackage[margin=0.1in]{geometry}
\usepackage{listings}
\usepackage{xcolor}
\usepackage{multicol}

% Java listings: same options as \lstdefinestyle{javaexam} in exam-test/main.tex
\lstset{
  language=Java,
  basicstyle=\ttfamily\footnotesize,
  keywordstyle=\bfseries\color{blue!55!black},
  commentstyle=\ttfamily\color{black!50},
  stringstyle=\color{green!50!black},
  columns=fullflexible,
  keepspaces=true,
  tabsize=2,
  showstringspaces=false,
  breaklines=true,
  breakatwhitespace=true,
  frame=single,
  rulecolor=\color{black!35},
  backgroundcolor=\color{black!4},
  xleftmargin=0.5em,
  xrightmargin=0.5em,
  aboveskip=4pt,
  belowskip=4pt,
}

\title{Java Coding Challenge Cheatsheet}
\author{Software Engineering Technical Reference}
\date{}

\begin{document}

% \maketitle

\begin{multicols}{2}




\section{Collections: List, Set, Map, PQ}

\subsection{List (ArrayList)}
\begin{lstlisting}
// mutable
List<Integer> list = new ArrayList<>(List.of(1, 2, 3));
list.add(4);                          // O(1) amort.; [1,2,3,4]
list.remove(list.size() - 1);         // O(1); [1,2,3]
int head = list.get(0);               // O(1); head == 1
list.set(0, 5);                       // O(1); [5,2,3]
list.contains(5);                     // O(n); true
list.isEmpty();                       // O(1); false
list.size();                          // O(1); 3
list.sort(Comparator.naturalOrder()); // O(n log n); [2,3,5]
\end{lstlisting}

\subsection{PriorityQueue (running min / max)}
\begin{lstlisting}
// min-heap (default Comparator)
PriorityQueue<Integer> minPq = new PriorityQueue<>();
minPq.offer(5);                   // O(log n)
minPq.offer(1);
minPq.offer(3);
int runningMin = minPq.peek();    // O(1); 1
minPq.poll();                     // O(log n); removes returns 1

// max-heap
Comparator<Integer> reverse = Comparator.reverseOrder();
PriorityQueue<Integer> maxPq = new PriorityQueue<>(reverse);
maxPq.offer(5);
maxPq.offer(1);
maxPq.offer(3);
int runningMax = maxPq.peek();    // O(1); 5
\end{lstlisting}

\subsection{Set (HashSet / TreeSet)}
\begin{lstlisting}
Set<Integer> set = new HashSet<>();
set.add(1);                       // O(1); {1}
set.remove(1);                    // O(1); {}
set.contains(1);                  // O(1); false

// TreeSet<Integer> (sorted unique)
TreeSet<Integer> tSet = new TreeSet<>(List.of(1, 3, 5, 10));
int lo = tSet.first();            // O(log n); 1
int hi = tSet.last();             // O(log n); 10
Integer sup = tSet.higher(5);     // least upper bound of 5 -> 10
Integer inf = tSet.lower(5);      // greatest lower bound of 5 -> 3
\end{lstlisting}

\subsection{Map (HashMap)}
\begin{lstlisting}
Map<Integer, String> map = new HashMap<>();
map.put(1, "A");                  // O(1); {1->"A"}
String v1 = map.get(1);           // O(1); "A"
String v2 = map.getOrDefault(2, "N/A"); // O(1); "N/A"
map.containsKey(1);               // O(1); true
map.putIfAbsent(1, "B");          // O(1); leaves "A"
map.remove(1);                    // O(1); {}

map.put(2, "X");
Set<Integer> keys = map.keySet();       // {2};
Collection<String> vals = map.values(); // {"X"};
for (Integer k : map.keySet()) {
  String s = map.get(k);                // "X" when k == 2
}
for (String v : map.values()) {
  boolean ok = v.startsWith("X");       // true
}
\end{lstlisting}




\columnbreak
\section{Comparator (lambda)}

\begin{lstlisting}
// Comparator<Integer>
Comparator<Integer> desc = (a, b) -> Integer.compare(b, a);
Comparator<Integer> byAbs = Comparator.comparingInt(Math::abs);
Comparator<Integer> natural = Comparator.naturalOrder();
Comparator<Integer> reverse = Comparator.reverseOrder();

int[] arr = {3, -1, -4, 2};
Arrays.sort(arr); // arr == {-4, -1, 2, 3} ascending

arr = Arrays.stream(arr)
            .boxed() // IntStream -> Stream<Integer>
            .sorted(byAbs)
            .mapToInt(Integer::intValue)
            .toArray();
// by absolute value: {-1, 2, 3, -4}
\end{lstlisting}

\section{Type Conversions}

\subsection{List/Set to Array}
\begin{lstlisting}
// List<Integer> -> int[]
List<Integer> nums = new ArrayList<>(List.of(10, 20, 30));
// mutable (deep copy semantics)
int[] prim = nums.stream()
                 .mapToInt(i -> i) // primitive returns deep copy
                 .toArray();

// List<String> -> String[]
List<String> words = new ArrayList<>(List.of("hi", "bye"));
// mutable (shallow copy semantics)
String[] arr = words.toArray(new String[0]);
// mutable (shallow copy semantics)
String[] array = words.stream()
                      .toArray(String[]::new);
// arr and array both {"hi","bye"}; different array objects
\end{lstlisting}

\subsection{Array to List/Set}
\begin{lstlisting}
// int[] -> List / Set
int[] pArr = {1, 2, 2, 3};
// mutable (deep copy semantics)
List<Integer> l4 = Arrays.stream(pArr)
                         .boxed() // IntStream -> Stream<Integer>
                         .collect(Collectors.toList());
// mutable (deep copy semantics)
Set<Integer> sInt = Arrays.stream(pArr)
                          .boxed()
                          .collect(Collectors.toSet());

// Object[] -> List
String[] arr = {"A", "B"};
// mutable wrapper (move semantics)
List<String> l1 = Arrays.asList(arr); arr = null;
// mutable (shallow copy semantics)
List<String> l2 = new ArrayList<>(Arrays.asList(arr));
// immutable, read-only (shallow copy semantics)
List<String> l3 = List.of(arr);

// Object[] -> Set (unique)
String[] arrDup = {"A", "A", "B"};
Set<String> sObj = new HashSet<>(Arrays.asList(arrDup));
// sObj == {"A", "B"} (order not fixed)
\end{lstlisting}

\subsection{List and Set}
\begin{lstlisting}
// List -> Set
List<Integer> sourceList = new ArrayList<>(List.of(1, 1, 2));
Set<Integer> targetSet = new HashSet<>(sourceList); // {1, 2}

// Set -> List
Set<Integer> sourceSet = new HashSet<>(List.of(3, 2));
List<Integer> targetList = new ArrayList<>(sourceSet); // {3, 2}
\end{lstlisting}





\columnbreak
\section{Array Initialization}

\subsection{Primitive Arrays}
\begin{lstlisting}
int[] arr = new int[3];          // {0, 0, 0}
Arrays.fill(arr, -1);            // {-1, -1, -1}

int[][] matrix = new int[2][2];
for (int[] row : matrix) {
  Arrays.fill(row, 7);
}
// matrix[0] == {7,7}; matrix[1] == {7,7}

// int[] literal
int[] staticArr = {1, 2, 3, 4, 5};
\end{lstlisting}

\section{Stream API Usage}

\subsection{Common Operations}
\begin{lstlisting}
// Stream<String> pipeline -> List<String>
List<String> list = new ArrayList<>(
    List.of("Alpha", "ax", "Beta", "Atlas"));
// mutable (copy semantics, terminal collect)
List<String> result = list.stream()
                          .filter(s -> s.startsWith("A")) // Alpha, Atlas
                          .map(String::toLowerCase)
                          .distinct()
                          .sorted()
                          .limit(5)
                          .collect(Collectors.toList());
// result == ["alpha", "atlas"]
\end{lstlisting}

\subsection{Primitive Streams (IntStream)}
\begin{lstlisting}
// IntStream reductions
int sum = IntStream.of(1, 2, 3).sum();              // 6
int max = IntStream.range(0, 10).max().getAsInt(); // 9 (end exclusive)
int min = IntStream.rangeClosed(0, 10).min().getAsInt(); // 0

// Stream<Integer> -> IntStream -> summary
List<Integer> nums = List.of(10, 20, 30);
IntSummaryStatistics stats = nums.stream()
                                 .mapToInt(i -> i) // Stream<Integer> -> IntStream
                                 .summaryStatistics();
// stats.getAverage() == 20.0; getCount() == 3
\end{lstlisting}

\subsection{Matching and Finding}
\begin{lstlisting}
// short-circuit terminal predicates
List<String> list = new ArrayList<>(List.of("", "abc", "ab"));
boolean any = list.stream()
                  .anyMatch(String::isEmpty);       // true ("")
boolean all = list.stream()
                  .allMatch(s -> s.length() > 2); // false ("ab")
Optional<String> first = list.stream()
                             .findFirst();         // Optional.of("")
\end{lstlisting}

\columnbreak


\section{Generic Methods}
\begin{lstlisting}
// type parameter T scoped to method
public static <T> void printList(List<T> list) {
  for (T item : list) {
    System.out.println(item);
  }
}
\end{lstlisting}

\section{Generic Classes}
\begin{lstlisting}
// type parameter T scoped to class
public class Box<T> {
  private T value;

  public void set(T value) {
    this.value = value;
  }
}
\end{lstlisting}


\section{String}
\begin{lstlisting}
// immutable
String raw = "  /api/v1/users  ";
String t = raw.trim();                  // "/api/v1/users"
boolean hit = t.startsWith("/api");     // true
String[] parts = t.split("/");          // {"", "api","v1","users"}
String seg = t.substring(5, 7);         // "v1" [i, j)
int at = t.indexOf("v1");               // 5
char c = t.charAt(at);                  // 'v'
"API".equalsIgnoreCase(parts[1]);       // true
int n = t.length();                     // 13
boolean empty = t.isEmpty();            // false
\end{lstlisting}

\subsection{String and \texttt{char[]}}
\begin{lstlisting}
String msg = "hello";
char[] buf = msg.toCharArray();         // mutable (copy semantics)
buf[0] = 'H';
String out = new String(buf);           // "Hello"
String share = String.valueOf(buf);     // "Hello"
buf[1] = 'i';                           // buf -> "Hillo"
                                        // out -> "Hello"
                                        // share -> "Hello"
\end{lstlisting}

\subsection{StringBuilder (mutable)}
\begin{lstlisting}
// mutable
StringBuilder sb = new StringBuilder("items=");
sb.append(10).append(',').append(20);   // O(1); "items=10,20"
sb.setLength(sb.length() - 1);          // O(1); drop trailing char
String csv = sb.toString();             // O(1); immutable snapshot
sb.append(',').append(30);              // O(1); "items=10,2,30"
String longer = sb.toString();          // O(1); "items=10,2,30"
int len = sb.length();                  // O(1); 13
\end{lstlisting}




\end{multicols}

\end{document}